% Independent preview; original geometry and numeric relationships retained.
% 原创教学示意；二叉树平均与两组节点之间的桥边。
\begin{tikzpicture}[x={\dimexpr\linewidth/169\relax},y=1mm,
  font=\figurefont\footnotesize,text=NNInk]
  \node[nn hidden,minimum size=6mm] (root) at (39,29) {$h$};
  \foreach \i/\xx in {1/18,2/60}{
    \node[nn hidden,minimum size=5mm] (a\i) at (\xx,13) {};
    \draw[nn operator path] (a\i)--(root);
  }
  \foreach \i/\xx/\parent in {1/8/1,2/28/1,3/50/2,4/70/2}{
    \node[nn hidden,minimum size=4mm] (b\i) at (\xx,-3) {};
    \draw[nn operator path] (b\i)--(a\parent);
  }
  \foreach \i/\xx/\parent in {1/3/1,2/13/1,3/23/2,4/33/2,5/45/3,6/55/3,7/65/4,8/75/4}{
    \node[nn input,minimum size=4mm] (x\i) at (\xx,-20) {};
    \draw[nn flow] (x\i)--(b\parent);
  }
  \node at (39,-32) {$h=\frac18\sum_{j=1}^{8}x_j$};
  \draw[nn link] (87,-38)--(87,35);
  \foreach \i/\xx/\yy in {1/103/17,2/103/-17,3/117/0,4/146/0,5/160/17,6/160/-17}{
    \ifnum\i=3
      \node[nn hidden,minimum size=6mm] (v\i) at (\xx,\yy) {};
    \else\ifnum\i=4
      \node[nn hidden,minimum size=6mm] (v\i) at (\xx,\yy) {};
    \else
      \node[nn input,minimum size=6mm] (v\i) at (\xx,\yy) {};
    \fi\fi
  }
  \foreach \a/\b in {1/2,1/3,2/3,4/5,4/6,5/6}{\draw[nn link] (v\a)--(v\b);}
  \draw[nn operator path,<->] (v3)--(v4);
  \node[nn heading] at (132,-32) {跨组桥边};
\end{tikzpicture}
